\begin{abstract}
Software maintainability is a critical quality attribute that directly affects system evolution cost and reliability. Improving maintainability often requires architectural-level modifications demanding specialized expertise and significant manual effort. While Large Language Models (LLMs) have shown strong capabilities in code understanding and generation, their potential for systematic code-level software quality improvement guided by architectural context remains largely unexplored.

This thesis investigates whether LLMs can support software maintainability improvement through code-level tactic implementation guided by architectural context. The work proceeds in three sequential empirical studies. The first study evaluates an end-to-end pipeline that detects repository architecture, selects an appropriate tactic from a predefined catalog, and implements it through LLM-generated code modifications, measuring the Maintainability Index (MI) before and after the intervention across 162 open-source Python backend repositories. The second study validates the architecture detection component independently by evaluating five LLMs across four context configurations on a manually labeled dataset of 57 repositories. The third study applies tactic implementation using validated architecture labels to isolate the effect of detection quality on implementation outcomes.

The end-to-end pipeline produces a statistically significant but modest overall MI improvement ($\overline{\Delta MI} = +0.48$, $p = 0.001$, $r = 0.28$), with 13.6\% of completed repositories improving and a pipeline failure rate of 25.3\%. Using validated architecture labels raises the improvement rate to 42.9\% and the mean gain to $+1.484$, demonstrating that architecture detection quality propagates into downstream tactic selection.

The architecture detection study shows that the best configuration — Qwen3-coder-next with file tree and import dependency graph — achieves 70.2\% accuracy and 0.65 macro-F1, exceeding the majority-class baseline by 12.3 percentage points. Import dependency graphs consistently improve classification across all five evaluated models, while code signatures degrade or do not improve performance. LLM confidence scores are poorly calibrated and provide no reliable signal for filtering uncertain predictions. All models systematically confuse modular monolith and layered architectures, a limitation that does not diminish with richer input. The multiple-comparison issue across 20 model-prompt pairs bounds the statistical significance of individual comparisons.

The strongest maintainability improvements occur in small, script-based repositories where the Decomposability tactic can extract cohesive modules from monolithic files with near-zero MI scores. Repository size is the dominant moderating factor: on large repositories, individual file-level changes are absorbed by averaging and produce negligible aggregate MI changes. The implementation stage consistently operates at the code level rather than the architectural level — package count remained unchanged across all repositories and fan-out changed in fewer than 25\% of cases — confirming the scale gap observed in prior LLM-based refactoring research. A detailed architectural analysis reveals that the Maintainability Index does not distinguish genuine architectural restructuring (e.g., Paper2Rebuttal) from mechanical file splitting (e.g., webapp-color), supporting the decision to frame the contributions as code-level rather than architecture-aware.

These findings indicate that LLMs can support code-level maintainability improvement as a targeted triage and transformation tool, particularly for structurally simple repositories with clear decomposition opportunities. Full autonomous architectural refactoring requires advances in architecture detection precision, implementation scope at the module boundary level, and behavioral verification.

\vspace{1em}
\noindent
\textbf{Keywords:} software maintainability, architectural tactics, large language models, automated software improvement, static analysis, software architecture.
\end{abstract}
